\documentclass[12pt,a4paper]{article}
\usepackage[utf8]{inputenc}
\usepackage{graphicx}   % For including images
\usepackage{float}  % Allows [H] to fix figure placement
\usepackage{booktabs}   % For nicer tables
\usepackage{geometry}
\geometry{margin=1in}
\usepackage{caption}

\title{Insurance Portfolio Risk Analysis}
\author{Lindokuhle Ncini}
\date{30 March 2026}

\begin{document}

\maketitle

\section{Introduction}

This project analyzes an insurance portfolio to determine its \textbf{risk exposure}. 
We calculated \textbf{expected losses, Value at Risk (VaR), Conditional VaR (CVaR)}, 
and conducted a \textbf{Monte Carlo simulation} to assess potential extreme losses. 
Graphs and summary tables illustrate key findings.

\section{Portfolio Overview}

The dataset contains multiple insurance policies with the following columns:

\begin{itemize}
    \item \texttt{policy\_id} -- unique policy identifier
    \item \texttt{exposure} -- monetary exposure of each policy
    \item \texttt{loss\_prob} -- probability of loss
    \item \texttt{losses} -- expected loss
\end{itemize}

\section{Summary Statistics}

\begin{table}[h!]
\centering
\begin{tabular}{lc}
\toprule
\textbf{Metric} & \textbf{Value} \\
\midrule
Total Expected Loss & 2,528,305.57 \\
Analytical 95\% VaR & 7,119.42 \\
Analytical 95\% CVaR & 8,572.66 \\
Monte Carlo Mean Loss & 2,529,849.40 \\
Monte Carlo 95\% VaR & 2,981,595.75 \\
Monte Carlo 99\% VaR & 3,177,618.88 \\
\bottomrule
\end{tabular}
\caption{Summary of key portfolio risk metrics.}
\end{table}

\section{Visual Analysis}

\subsection{Loss Distribution with VaR and CVaR}

The histogram shows how expected losses are distributed across policies. 
Red dashed line = 95\% VaR, Orange dashed line = 95\% CVaR.

\begin{figure}[H]
\centering
\includegraphics[width=0.8\textwidth]{loss_distribution.png}
\caption{Loss distribution with VaR and CVaR.}
\end{figure}

\subsection{Monte Carlo Simulation of Total Losses}

Shows 10,000 simulated portfolio outcomes. Red line = 95\% VaR, Orange line = 99\% VaR. 
Highlights extreme loss scenarios.

\begin{figure}[H]
\centering
\includegraphics[width=0.8\textwidth]{monte_carlo_losses.png}
\caption{Monte Carlo simulation of total portfolio losses.}
\end{figure}



\section{Key Findings}

\begin{itemize}
    \item Total Expected Loss gives the expected total loss for the portfolio.
    \item 95\% VaR indicates the loss threshold not expected to be exceeded in 95\% of cases.
    \item 95\% CVaR shows the average of the worst 5\% of losses, highlighting extreme risk.
    \item Monte Carlo simulation confirms analytical metrics and provides a visual understanding of extreme loss scenarios.
\end{itemize}

\section{Conclusion}

The portfolio exhibits predictable expected losses, but extreme scenarios can lead to substantially higher losses. 
Management should focus on high-exposure, high-probability policies to mitigate risk.

\newpage  % Forces references to appear on a separate page
\section*{References}

\begin{enumerate}
    \item Hull, J. (2021) \textit{Risk Management and Financial Institutions}, 5th edition. Hoboken: Wiley.
    \item Glasserman, P. (2003) \textit{Monte Carlo Methods in Financial Engineering}. New York: Springer.
    \item McNeil, R., Frey, R. and Embrechts, P. (2015) \textit{Quantitative Risk Management: Concepts, Techniques and Tools}, 2nd edition. Princeton: Princeton University Press.

    \item Basel Committee on Banking Supervision (2011) \textit{Principles for the Sound Management of Operational Risk}. Basel: Bank for International Settlements.
\end{enumerate}

\end{document}
